\documentclass[11pt]{article}
\usepackage[margin=1in]{geometry}
\usepackage{amsmath,amssymb,amsthm}
\usepackage[hidelinks]{hyperref}
\newtheorem{theorem}{Theorem}
\title{The exact zero-free gap for\\the spectrum \(\boldsymbol{\{\pm1,\pm4,\pm7\}}\)}
\author{}
\date{27 November 2026}
\begin{document}
\maketitle
\vspace{-2em}
\begin{abstract}
For \(\{\pm N,\pm4N,\pm7N\}\) the supremal zero-free
interval has length \(\arccos(\xi)/(\pi N)\), where
\(\xi=-0.3639080166\ldots\) is an explicitly isolated algebraic
number. The supremum is not attained. On the critical closed interval,
there is a unique nonnegative polynomial up to positive scale, and a
unique probability measure annihilating the spectrum. That measure
requires exactly six atoms.
\end{abstract}

\section{Statement and algebraic data}
For a symmetric set \(S\) of nonzero integers let \(M(S)\) be the
supremum of lengths of intervals in \(\mathbb R/\mathbb Z\) on which a
real trigonometric polynomial with spectrum contained in \(S\) has no
zeros. Intervals may equivalently be taken in the periodic lift.
Kozma and Oravecz ask for \(M(\{\pm(N+jb):0\le j\le K\})\) when
\(b\ge 2N\) \cite{KO}. The ProveIt spectral-filter chapter relates the
problem to minimum-width positive annihilators \cite{P}. We settle the
subfamily \(K=2,\ b=3N\), without using RH or GRH.

\begin{theorem}
For every positive integer \(N\),
\[
 M(\{\pm N,\pm4N,\pm7N\})=\frac{\arccos(\xi)}{\pi N}
       \simeq \frac{0.6185577534840}{N}.
\]
No zero-free interval has this length, even with both endpoints removed.
On a fixed closed interval of this length, the cone of nonnegative
polynomials on this spectrum is a single ray, and the probability
measure annihilating the three frequencies is unique and has six atoms.
\end{theorem}

Write \(T_j(\cos\phi)=\cos(j\phi)\) and set
\[
 U(x)=64x^6-112x^4+56x^2-7,\qquad V(x)=8x^3-8x+x^{-1}.
\]
Thus \(U=T_7/x\) away from zero and \(V=T_4/x\). Define
\(H(x,B)=B U(x)+V(x)\) and \(F=\partial_xH\).
The triple \((B_*,u,v)\) is the unique solution in the boxes
\begin{align*}
 J_B&=[0.9208242456520,\ 0.9208242456522],\\
 J_u&=[0.2680823592348,\ 0.2680823592351],\\
 J_v&=[0.8545996486673,\ 0.8545996486676]
\end{align*}
of \(F(u,B_*)=F(v,B_*)=0\) and \(H(v,B_*)=H(u,B_*)\).
Put \(A=-H(u,B_*)\) and \(P(x)=T_4(x)+Ax+B_*T_7(x)\).
The number \(\xi\) is the unique root of \(P\) in
\((-0.363908019,-0.363908014)\).

\section{Isolation and sign of the polynomial}
The reproducible sign checks in Section~\ref{sec:check} give
\(\partial_x F>0\) on each of \(J_u\times J_B\) and \(J_v\times J_B\),
and \(F<0\) at the left and \(F>0\) at the right endpoint of each
\(x\)-box, uniformly in \(B\). Hence there are unique differentiable
root branches \(u(B),v(B)\). The derivative of
\(\Delta(B)=H(v(B),B)-H(u(B),B)\) is \(U(v(B))-U(u(B))>0\);
the implicit-derivative contributions vanish by stationarity.
The endpoints of \(J_B\) give \(\Delta(B_{\rm left})<0<
\Delta(B_{\rm right})\), proving existence and uniqueness. Clearing the
nonzero denominators gives a polynomial system over \(\mathbb Q\).
The solution has algebraic coordinates by quantifier elimination for
real closed fields: the same unique solution exists over the real
algebraic numbers.

Since \(P=x(A+H(x,B_*))\) away from zero, \(u,v\) are double zeros:
\[
 P(x)=(x-u)^2(x-v)^2G(x),\qquad
 G(x)=c_3x^3+c_2x^2+c_1x+c_0 .
\]
With \(h=u+v,\ k=uv,\ w=h^2+2k\), comparing the top four coefficients gives
\begin{align*}
 c_3&=64B_*, & c_2&=2hc_3,\\
 c_1&=-112B_*+2hc_2-wc_3, &
 c_0&=8+2hc_1-wc_2+2hkc_3.
\end{align*}
Let \(x_-=-0.363908019,\ x_+=-0.363908014\). The exact checks give
\(c_3>0\), \(G'(x_-)>0\), \(G''(x_-)>0\), and \(G(x_-)<0<G(x_+)\).
Thus \(G\) is strictly increasing on \([x_-,1]\), with unique root
\(\xi\in(x_-,x_+)\). Therefore \(P\) is nonnegative on \([\xi,1]\)
and vanishes there exactly at \(\xi,u,v\).
The same checks give \(U(u)<U(\xi)<U(v)\).

\section{Sharp dual certificate and approximation}
It is enough to work at \(N=1\): the changes of variable \(s=Nt\) and
\(t=s/N\) biject the classes of lifted polynomials and scale interval
lengths by \(N\) and \(1/N\), respectively. Write \(L=\arccos(\xi)/\pi\).
For \(x\in\{\xi,u,v\}\) set \(t_x=\arccos(x)/(2\pi)\), and define
\[
 \theta=\frac{U(v)-U(\xi)}{U(v)-U(u)},\qquad
 w_\xi=1,\quad w_u=\frac{|\xi|\theta}{u},\quad
 w_v=\frac{|\xi|(1-\theta)}{v}.
\]
The weights are positive and satisfy
\(\sum_x w_xx=\sum_xw_xT_7(x)=0\); summing \(P(x)=0\) gives
\(\sum_xw_xT_4(x)=0\). Split each weight equally between \(t_x\) and
\(-t_x\), and divide by the total mass. The resulting measure \(\mu\)
annihilates all six sine and cosine modes of frequencies \(1,4,7\),
since the sine modes cancel in pairs. Its support lies in
\([-L/2,L/2]\), with four atoms strictly inside.

If a polynomial is zero-free on an interval of length at least \(L\),
center a subinterval of length \(L\) at zero by translation and change
the sign of the polynomial if necessary. It is strictly positive
in the interior and nonnegative at the endpoints. Integration against
\(\mu\) is strictly positive, contradicting annihilation.
Conversely, for any \(\xi<a<0\), \(P(x)+\varepsilon x\) is strictly
positive on \([a,1]\) for sufficiently small \(\varepsilon>0\).
On \([a,0]\) the polynomial \(P\) has a positive minimum, and at
positive \(x\) the perturbation is positive. Substituting
\(x=\cos(2\pi t)\) yields zero-free intervals of lengths
\(\arccos(a)/\pi\) approaching \(L\) with no new frequencies.

\clearpage
\section{Uniqueness on the critical interval}
Let \(g\) be nonnegative on \([-L/2,L/2]\) with the allowed spectrum.
Its integral against \(\mu\) is zero, so \(g(\pm t_x)=0\) for each \(x\).
The even part is \(C_1T_1+C_4T_4+C_7T_7\) in \(x\). Evaluation at
\(u,v\) has rank two: the \(T_1,T_7\) submatrix has determinant
\(uv(U(v)-U(u))\ne 0\). As \(P\) vanishes at both points, the even
part must be a scalar multiple of \(P\).

For the odd part use
\(\sin(j\phi)=\sin(\phi)W_{j-1}(\cos\phi)\) with
\[
 W_0=1,\qquad W_3=8x^3-4x,\qquad
 W_6=64x^6-80x^4+24x^2-1 .
\]
The determinant with rows \((1,W_3(x),W_6(x))\) at \(x=u,v,\xi\)
lies strictly between \(4.54\) and \(4.55\) by the rational check.
The factors \(\sin(2\pi t_x)\) are nonzero; hence the odd part vanishes.
The scalar multiplying \(P\) is nonnegative because \(P(1)>0\).

For any annihilating probability measure on the interval,
integration of \(P(\cos(2\pi t))\) forces it to be supported
on the six zeros. The sine equations, by the same invertible matrix,
force the masses of each pair to agree. The cosine equations for
\(T_1,T_7\) have rank two in the three pair totals, so normalization
determines the pair totals uniquely. This is precisely \(\mu\);
all six masses are necessary.

\section{Exact replay}\label{sec:check}
Run \texttt{node certify.js} and \texttt{python3 verify.py}. Neither
program requires third-party packages. The first uses outward-rounded
integer intervals on a \(10^{-24}\) grid to check the original boxes;
only its printed summary uses floating point. The second uses
exact rational endpoints, isolates the stationary roots at the two
endpoints of a tighter \(B\)-box by bisection, and checks the cubic and
sine determinant on tighter boxes. These lie inside the original boxes,
so their isolated solution is the same. In particular,
\[
 -0.36390801664214987744459<\xi<-0.36390801664214987744456.
\]
The proof uses the exact inequalities rather than the displayed
approximation to \(L\). General finite off-resonance progressions are
not resolved here.

\begin{thebibliography}{9}
\bibitem{KO}
G. Kozma and F. Oravecz,
\emph{On the gaps between zeros of trigonometric polynomials},
Real Anal. Exchange 28 (2002/03), no.~2, 447--454.
\href{https://arxiv.org/abs/math/0601021}{arXiv:math/0601021},
Question~1.
\bibitem{P}
ProveIt,
\href{https://github.com/VladimirReshetnikov/ProveIt/blob/531fa8ce559ed7a34b1a962bd325472171741000/Analysis/GeneralizedRiemannHypothesis/Article/sections/fil.tex}{GRH article, spectral-filter chapter},
revision \texttt{531fa8ce559e}, paragraph on zero-free intervals
and the convolution bound.
\end{thebibliography}
\end{document}
